\documentclass[12pt,a4paper]{article}

\usepackage[utf8]{inputenc}
\usepackage[T1]{fontenc}
\usepackage[english]{babel}
\usepackage{geometry}
\geometry{margin=2.5cm}

\usepackage{amsmath,amssymb}
\usepackage{graphicx}
\usepackage{subcaption}
\usepackage{booktabs}
\usepackage{hyperref}
\usepackage{xcolor}
\usepackage{float}
\usepackage{enumitem}
\usepackage[expansion=false]{microtype}
\usepackage{parskip}
\usepackage{titlesec}
\usepackage{fancyhdr}
\setlength{\headheight}{15pt}
\usepackage{caption}
\usepackage{url}

\graphicspath{{../outputs/}}

\pagestyle{fancy}
\fancyhf{}
\rhead{Deep Learning -- Project III}
\lhead{Cat Image Synthesis}
\cfoot{\thepage}

\title{\textbf{Generative Models for Cat Image Synthesis}\\
       \large Deep Learning -- Project III}
\author{Anna Ostrowska}
\date{June 2026}

\begin{document}
\maketitle

\begin{abstract}
This report presents the implementation and evaluation of three generative
modelling approaches for synthesising cat images at $64\times64$ resolution:
a Deep Convolutional GAN trained with Wasserstein loss and gradient penalty
(DCGAN/WGAN-GP), a Vector-Quantized Variational Autoencoder with an
autoregressive PixelCNN prior (VQ-VAE), and a Denoising Diffusion Probabilistic
Model with DDIM fast sampling (DDPM).  All models are trained from scratch on
the CAT Dataset ($9\,997$ images) using a single consumer GPU with 6\,GiB VRAM.
Models are compared quantitatively using the Fr\'{e}chet Inception Distance (FID)
computed on 5\,000 generated images and qualitatively through visual inspection
and latent space interpolation.  DDPM achieves the lowest FID (24.36 at 50 DDIM
steps with EMA weights), substantially outperforming VQ-VAE (187.54) and
DCGAN (221.51).  A post-hoc analysis of the DCGAN training configuration reveals
that generator starvation caused by $n_{\text{critic}}=5$ likely contributed to
the high FID; corrected configurations reduce FID to 98.99 (standard DCGAN) and
158.48 (WGAN-GP $n_c=2$).  Ablation studies investigate the effect of DDIM
inference steps and EMA weights on DDPM quality, and the sensitivity of DCGAN to
latent dimension and learning rate.  An exploratory extension trains DCGAN on a
combined cats-and-dogs dataset, finding a suggestive improvement in cat FID that
warrants further investigation.  Source code and model checkpoints are provided
for full reproducibility.
\end{abstract}

\tableofcontents
\newpage

\section{Introduction}
\label{sec:intro}

This report presents the design, training, and evaluation of three classes of
generative image models applied to a dataset of cat photographs.  The three
approaches investigated are:

\begin{enumerate}[noitemsep]
  \item \textbf{DCGAN} with Wasserstein loss and gradient penalty (WGAN-GP),
  \item \textbf{VQ-VAE} with an autoregressive PixelCNN prior,
  \item \textbf{DDPM} (Denoising Diffusion Probabilistic Model) with DDIM
        sampling.
\end{enumerate}

Models are compared quantitatively using the Fr\'{e}chet Inception Distance (FID)
and qualitatively through visual inspection of generated samples and latent
interpolation strips.  Additionally, an ablation study is conducted for each
model class, and a supplementary experiment trains a DCGAN on a combined
cats-and-dogs dataset.

\paragraph{Deviations from the project plan.}
Three aspects of the implementation diverged from the original project plan
(KM0, submitted 19 May 2026):

\begin{enumerate}[noitemsep]
  \item \textbf{Label smoothing.}
        The plan stated ``WGAN-GP loss with one-sided label smoothing
        (real labels $= 0.9$) is used for DCGAN.''  In practice, this
        parameter -- present in the shared codebase for standard BCE-GAN
        variants -- operates on binary cross-entropy targets and has no
        effect on the WGAN-GP loss path, which uses an unbounded scalar
        critic score rather than class probabilities.  Label smoothing was
        therefore inactive during WGAN-GP training.

  \item \textbf{Latent interpolation.}
        The plan described linear interpolation between latent vectors.
        We switched to spherical linear interpolation (slerp) for models
        whose prior is $\mathcal{N}(0,I)$ (DCGAN, DDPM), because linear
        interpolation crosses a low-density region near the origin and
        produces systematically blurry intermediate frames
        (see Section~\ref{sec:interpolation}).  Linear interpolation was
        retained for VQ-VAE, where the prior operates over discrete tokens
        rather than a Gaussian.

  \item \textbf{DCGAN ablation depth.}
        The plan specified 50-epoch ablation runs.  In practice, ablations
        were extended to 100 epochs and repeated across three random seeds
        to reduce the risk of comparing under-trained models and to quantify
        seed variance, reported as mean $\pm$ population standard deviation
        in Table~\ref{tab:dcgan_ablation}.
\end{enumerate}

\section{Background and Related Work}
\label{sec:background}

\subsection{Generative Adversarial Networks}

Generative Adversarial Networks (GANs) \cite{gan} frame image synthesis as a
two-player game between a Generator $G$ and a Discriminator $D$.  The generator
maps samples $\mathbf{z} \sim p(\mathbf{z})$ from a simple prior (typically
$\mathcal{N}(0,I)$) to synthetic images, while the discriminator attempts to
distinguish real from generated data.  Training minimises the standard
min-max objective:
\begin{equation}
  \min_G \max_D\; \mathbb{E}_{x\sim p_r}[\log D(x)]
                + \mathbb{E}_{\mathbf{z}\sim p(\mathbf{z})}[\log(1-D(G(\mathbf{z})))].
\end{equation}
GANs are prone to \emph{mode collapse} -- a failure mode where the generator
produces only a small subset of the data distribution -- and to unstable
gradient dynamics.  The Deep Convolutional GAN (DCGAN) \cite{dcgan} introduced
architectural guidelines (all-convolutional networks, batch normalisation,
LeakyReLU activations) that substantially improved training stability.

\paragraph{WGAN and WGAN-GP.}
Arjovsky et al.\ \cite{wgan} proposed replacing the Jensen-Shannon divergence
objective with the Wasserstein-1 distance, which provides smoother gradients
and alleviates mode collapse.  However, their weight-clipping approach
introduces gradient pathologies.  Gulrajani et al.\ \cite{wgangp} resolved
this by enforcing the Lipschitz constraint via a gradient penalty term on
interpolated samples, yielding the Wasserstein GAN with Gradient Penalty
(WGAN-GP).  Crucially, the WGAN-GP critic outputs an unbounded scalar score
rather than a probability, so binary label targets such as $\{0, 1\}$ or
their smoothed variants are not used in the Wasserstein loss path.

\subsection{Variational Autoencoders and VQ-VAE}

Variational Autoencoders (VAEs) \cite{vae} train an encoder--decoder pair by
maximising the evidence lower bound (ELBO):
\begin{equation}
  \mathcal{L}_{\text{VAE}} = \mathbb{E}_{q_\phi(z|x)}[\log p_\theta(x|z)]
                            - \mathrm{KL}(q_\phi(z|x) \| p(z)),
\end{equation}
where $q_\phi$ is the approximate posterior and $p_\theta$ the decoder.
The continuous latent space allows smooth interpolation but tends to produce
blurry reconstructions.

The Vector-Quantized VAE (VQ-VAE) \cite{vqvae} replaces the continuous latent
code with a discrete codebook.  The encoder output is replaced by the nearest
codebook entry (straight-through gradient); codebook updates use exponential
moving averages (EMA) for stability.  Because the prior $p(z)$ is non-trivial
(discrete, structured), a separate autoregressive model (PixelCNN \cite{pixelcnn})
is trained over the quantised token sequences after the autoencoder converges.

\subsection{Denoising Diffusion Probabilistic Models}

DDPM \cite{ddpm} defines a \emph{forward process} that gradually adds Gaussian
noise to data:
\begin{equation}
  q(x_t | x_{t-1}) = \mathcal{N}(x_t;\, \sqrt{1-\beta_t}\, x_{t-1},\; \beta_t I),
\end{equation}
and a learnable \emph{reverse process} $p_\theta(x_{t-1}|x_t)$ implemented as
a U-Net that predicts the noise $\epsilon$ added at each step.  The training
objective is:
\begin{equation}
  \mathcal{L}_{\text{DDPM}} = \mathbb{E}_{t,x_0,\epsilon}
    \left[\|\epsilon - \epsilon_\theta(x_t, t)\|^2\right].
\end{equation}
Dhariwal and Nichol \cite{diffusionbeatsgan} demonstrated that diffusion models
with classifier guidance surpass GANs on FID benchmarks.

\paragraph{DDIM.}
Song et al.\ \cite{ddim} showed that the same trained model can be used with a
deterministic, non-Markovian sampling process (Denoising Diffusion Implicit
Models) requiring as few as 20--50 function evaluations rather than the full
$T=1000$ steps, at negligible quality loss.

\subsection{Fr\'{e}chet Inception Distance}

FID \cite{fid} measures the Fr\'{e}chet (Wasserstein-2) distance between the
distributions of real and generated images in the feature space of an
Inception-v3 network:
\begin{equation}
  \mathrm{FID} = \|\mu_r - \mu_g\|^2
               + \mathrm{tr}(\Sigma_r + \Sigma_g - 2(\Sigma_r \Sigma_g)^{1/2}),
\end{equation}
where $(\mu_r, \Sigma_r)$ and $(\mu_g, \Sigma_g)$ are the mean and covariance
of the Inception activations for real and generated sets.  Lower FID indicates
closer distributional similarity.  We use \texttt{torch-fidelity} with 5\,000
real and 5\,000 generated images.

\section{Dataset and Preprocessing}
\label{sec:data}

\subsection{Cat Dataset}
The primary dataset is the CAT Dataset \cite{catdataset}, consisting of
$9\,997$ RGB cat photographs with heterogeneous resolutions, organised in
subdirectories \texttt{CAT\_00} through \texttt{CAT\_06}.

\subsection{Preprocessing Pipeline}
A uniform preprocessing pipeline was applied to all images:

\begin{enumerate}[noitemsep]
  \item \textbf{Resize} the shorter edge to $64$ pixels.
  \item \textbf{CenterCrop} to $64 \times 64$.
  \item \textbf{Normalize} to $[-1, 1]$ via
        $x \leftarrow (x - 0.5) / 0.5$ per channel.
\end{enumerate}

All models operate on $64 \times 64 \times 3$ tensors.
Figure~\ref{fig:data_samples} shows a random selection of preprocessed images.
A fixed random seed (42) is used for all experiments.  DataLoaders use
\texttt{drop\_last=True}, giving $\lfloor 9997 / 64 \rfloor = 156$ batches
per epoch.

\begin{figure}[H]
  \centering
  \includegraphics[width=0.7\linewidth]{data_samples_64.png}
  \caption{Random sample of real cat images after $64\times64$ center-crop preprocessing.}
  \label{fig:data_samples}
\end{figure}

\subsection{Data Augmentation}
An optional augmentation pipeline was implemented.  When enabled, the following
transforms are applied before normalization:

\begin{itemize}[noitemsep]
  \item \texttt{RandomResizedCrop} (scale $0.85$--$1.0$),
  \item \texttt{RandomHorizontalFlip} ($p=0.5$),
  \item \texttt{ColorJitter} (brightness/contrast/saturation $= 0.2$, hue $= 0.05$),
  \item \texttt{RandomGrayscale} ($p=0.05$),
  \item \texttt{RandomRotation} ($\pm10^{\circ}$),
  \item \texttt{RandomErasing} ($p=0.3$, scale $2$--$8\%$).
\end{itemize}

The baseline models are trained without augmentation; augmentation effects are
investigated in the ablation study (Section~\ref{sec:ablation}).

\section{Methods}
\label{sec:methods}

\subsection{DCGAN with WGAN-GP}
\label{sec:dcgan}

\paragraph{Architecture.}
The Generator $G$ maps a latent vector $\mathbf{z} \sim \mathcal{N}(0,I)$,
$\mathbf{z} \in \mathbb{R}^{128}$, to a $64\times64\times3$ image through four
transposed-convolution blocks (BatchNorm + ReLU), ending in a Tanh activation.
The Discriminator $D$ mirrors this with four strided-convolution blocks
(LeakyReLU, $\alpha=0.2$) and outputs a scalar logit without sigmoid.
The architecture follows DCGAN guidelines \cite{dcgan}.

\textbf{Implementation note.}
The discriminator retains Batch Normalisation layers from the DCGAN architecture,
which is a known sub-optimal choice for WGAN-GP: BatchNorm mixes statistics
across samples in a batch, which interferes with the per-sample gradient penalty
computation.  The original WGAN-GP paper recommends Layer Normalisation in the
critic.  We did not modify the critic normalisation; this may have slightly
reduced training stability and is a limitation of the current implementation.

\paragraph{Training objective.}
We adopt the Wasserstein loss with gradient penalty (WGAN-GP \cite{wgangp}):
\begin{equation}
  \mathcal{L}_D = \mathbb{E}_{\tilde{x}\sim P_g}[D(\tilde{x})]
                - \mathbb{E}_{x\sim P_r}[D(x)]
                + \lambda \, \mathbb{E}_{\hat{x}}
                  \!\left[\left(\|\nabla_{\hat{x}} D(\hat{x})\|_2 - 1\right)^2\right],
\end{equation}
with $\lambda = 10$ and $\hat{x}$ sampled uniformly along straight lines between
real and generated data.  The discriminator (critic) is updated
$n_{\text{critic}} = 5$ times per generator step.  Because the WGAN-GP critic
returns an unbounded scalar rather than a class probability, binary label targets
-- and any smoothing thereof -- have no role in this loss formulation and were
not applied.

\paragraph{Hyperparameters.}
Trained for 100 epochs, batch size 64, Adam ($\beta_1=0.5$, $\beta_2=0.999$),
learning rate $2 \times 10^{-4}$ for both $G$ and $D$.

\subsection{VQ-VAE + PixelCNN Prior}
\label{sec:vqvae}

\paragraph{Architecture.}
The VQ-VAE \cite{vqvae} consists of an encoder $E$ that maps a $64\times64$
image to a $16\times16$ discrete codebook index map, a vector-quantization layer
with codebook size $K=512$ and embedding dimension $d=64$, and a decoder $D$
that reconstructs the image from quantized embeddings.  The encoder and decoder
each use strided-convolution and transposed-convolution blocks with two residual
layers.  Codebook updates use exponential moving averages (EMA) with decay
$\gamma = 0.99$.

Two-phase training:
\begin{enumerate}[noitemsep]
  \item \textbf{Phase 1 -- Autoencoder} (100 epochs): minimizes reconstruction
        loss + VQ commitment loss $\beta \|z_e - \text{sg}(z_q)\|^2$, $\beta=0.25$.
  \item \textbf{Phase 2 -- Prior} (100 epochs): a PixelCNN \cite{pixelcnn} with
        15 masked residual layers (128 filters each) is trained autoregressively
        on the flattened $16 \times 16 = 256$ token sequences from the frozen encoder.
\end{enumerate}

\paragraph{Sampling.}
At inference, the PixelCNN samples a $256$-token sequence autoregressively;
tokens are mapped to embeddings via the codebook and decoded by $D$.

\subsection{DDPM with DDIM Sampling}
\label{sec:ddpm}

\paragraph{Architecture.}
We implement the DDPM framework \cite{ddpm} with a U-Net backbone.  The U-Net
has channel multipliers $(1, 2, 2, 2)$, two residual blocks per resolution, and
self-attention at the $16\times16$ feature map level.  An exponential moving
average (EMA, decay $0.9999$) of model weights is maintained and used at
inference.

\paragraph{Diffusion schedule.}
Linear noise schedule with $T=1000$ steps, $\beta_1 = 10^{-4}$,
$\beta_T = 0.02$.  The forward process is:
\begin{equation}
  q(x_t | x_0) = \mathcal{N}(x_t;\, \sqrt{\bar\alpha_t}\, x_0,\;
                 (1 - \bar\alpha_t) I),
\end{equation}
where $\bar\alpha_t = \prod_{s=1}^t (1 - \beta_s)$.  The model predicts the
added noise $\epsilon_\theta(x_t, t)$ and is trained with simple $\ell_2$ loss.

\paragraph{DDIM Sampling.}
Deterministic DDIM \cite{ddim} allows generating samples with far fewer
function evaluations than the full $T=1000$ reverse chain.  At inference we
use 50 DDIM steps ($\eta = 0.0$) by default, giving a $20\times$ speedup.

\paragraph{Hyperparameters.}
Trained for 200 epochs, batch size 16 (limited by 6~GiB VRAM), Adam
($\beta_1=0.9$, $\beta_2=0.999$), learning rate $2 \times 10^{-4}$ with linear
warmup.

\section{Computational Considerations}
\label{sec:compute}

All experiments were run on a single NVIDIA GPU with 6~GiB VRAM and a Windows
host.  Key adaptations:

\begin{itemize}[noitemsep]
  \item \texttt{num\_workers=0} in all DataLoaders (Windows multiprocessing
        limitation).
  \item DDPM batch size reduced to 16 to fit within VRAM.
  \item DDIM sampling used instead of full reverse diffusion to cut inference
        time from $\approx8$\,h (1000 steps) to $\approx$25\,min (50 steps).
  \item VQ-VAE EMA codebook uses \texttt{scatter\_add\_} + \texttt{bincount}
        instead of \texttt{F.one\_hot} to avoid out-of-memory issues.
  \item Global seed 42 for reproducibility.
\end{itemize}

\section{Results}
\label{sec:results}

\subsection{FID Comparison}
\label{sec:fid}

Fr\'{e}chet Inception Distance (FID) \cite{fid} measures the distance between the
distribution of real images and generated images in the feature space of an
Inception-v3 network.  We generate 5\,000 samples per model and compare against
5\,000 real images drawn from the training set.

\begin{table}[H]
  \centering
  \caption{FID scores on the cat dataset ($64\times64$, 5\,000 samples).
           All results are single-run.  The DCGAN entry uses the original
           training configuration later found to be under-trained (see
           Section~\ref{sec:dcgan_bottleneck}); a corrected standard DCGAN
           achieves FID 98.99.}
  \label{tab:fid}
  \begin{tabular}{lcc}
    \toprule
    Model & FID $\downarrow$ & Notes \\
    \midrule
    DCGAN (WGAN-GP, $n_c{=}5$) & 221.51 & 100 ep., 3\,100 G-steps; under-trained$^\dagger$ \\
    DCGAN (standard, corrected) & 98.99  & 150 ep., 23\,400 G-steps; see \S\ref{sec:dcgan_bottleneck} \\
    VQ-VAE + PixelCNN  & 187.54 & 100 + 100 epochs \\
    DDPM (50 DDIM steps, EMA) & \textbf{24.36} & 200 epochs \\
    \bottomrule
  \end{tabular}
  \begin{flushleft}
  \small $^\dagger$ With $n_c{=}5$ the generator receives only 3\,100 weight
  updates over 100 epochs vs.\ 15\,600 for a standard 1:1 update ratio.
  \end{flushleft}
\end{table}

DDPM achieves a substantially lower FID, consistent with the literature
showing diffusion models outperforming GANs and VAEs on small-scale datasets
\cite{diffusionbeatsgan}.  The DCGAN result depends strongly on training
configuration: the original WGAN-GP setting ($n_c=5$) results in only
3\,100 generator updates over 100 epochs, which we show in
Section~\ref{sec:dcgan_bottleneck} is a likely source of under-training.
A corrected standard DCGAN achieves FID 98.99, substantially narrowing the
gap with VQ-VAE.

\textbf{Scope of this comparison.}
Two caveats apply.  First, the comparison is not compute-normalised: DDPM was
trained for $\approx$12\,h, DCGAN for $\approx$2\,h, and VQ-VAE for $\approx$4\,h.
The table reflects the best feasible configuration for each model under the
6\,GiB VRAM constraint, not a strictly equal-budget comparison.  Second, FID was
computed against a subset of the \emph{training} distribution; no held-out test
set was used.  The metric therefore measures similarity to available training
data rather than generalisation to unseen cat images.  Given that the main
comparisons are single-run, the large gap between DDPM and the other models is
reliably qualitative, while smaller differences should not be over-interpreted.

\subsection{Generated Samples}
\label{sec:samples}

\begin{figure}[H]
  \centering
  \begin{subfigure}[b]{0.32\linewidth}
    \includegraphics[width=\linewidth]{dcgan/samples_epoch0100.png}
    \caption{DCGAN (epoch 100)}
  \end{subfigure}
  \hfill
  \begin{subfigure}[b]{0.32\linewidth}
    \includegraphics[width=\linewidth]{vqvae/prior_samples.png}
    \caption{VQ-VAE + PixelCNN}
  \end{subfigure}
  \hfill
  \begin{subfigure}[b]{0.32\linewidth}
    \includegraphics[width=\linewidth]{ddpm/samples_epoch0200.png}
    \caption{DDPM (epoch 200)}
  \end{subfigure}
  \caption{Generated samples from each model at the end of training.}
  \label{fig:samples}
\end{figure}

DDPM samples show the most coherent cat structure with natural fur textures.
DCGAN samples exhibit some diversity but with noticeable artifacts and
occasional mode collapse symptoms (similar face structures).  VQ-VAE samples
are globally plausible but blurrier, which likely reflects limitations of the
PixelCNN prior rather than the autoencoder itself (see Section~\ref{sec:vqvae_recon}).

\subsection{Training Dynamics}
\label{sec:training_curves}

\begin{figure}[H]
  \centering
  \begin{subfigure}[b]{0.32\linewidth}
    \includegraphics[width=\linewidth]{dcgan/losses.png}
    \caption{DCGAN losses (G and D)}
  \end{subfigure}
  \hfill
  \begin{subfigure}[b]{0.32\linewidth}
    \includegraphics[width=\linewidth]{vqvae/losses.png}
    \caption{VQ-VAE reconstruction loss}
  \end{subfigure}
  \hfill
  \begin{subfigure}[b]{0.32\linewidth}
    \includegraphics[width=\linewidth]{ddpm/losses.png}
    \caption{DDPM noise-prediction loss}
  \end{subfigure}
  \caption{Training loss curves for all three models.}
  \label{fig:losses}
\end{figure}

\begin{figure}[H]
  \centering
  \includegraphics[width=\linewidth]{vqvae/prior_losses.png}
  \caption{VQ-VAE PixelCNN prior training loss (Phase 2, 100 epochs).}
  \label{fig:prior_losses}
\end{figure}

DCGAN training shows the characteristic adversarial oscillation: generator and
discriminator losses compete without a stable equilibrium.  VQ-VAE reconstruction
loss decreases steadily throughout Phase 1 and the PixelCNN prior loss decreases
more slowly in Phase 2.  DDPM loss decreases monotonically and smoothly,
reflecting the well-defined mean-squared-error noise-prediction objective.

All models were trained on the full dataset without a validation split;
early stopping was therefore not applied.  Sample grids saved at regular
checkpoints (shown in Section~\ref{sec:progression}) serve as the primary
qualitative proxy for monitoring training progress and detecting mode collapse.
Per-epoch FID evaluation was not feasible due to its computational cost
($\approx$5\,min per evaluation).

\subsection{Sample Progression}
\label{sec:progression}

\begin{figure}[H]
  \centering
  \begin{subfigure}[b]{0.32\linewidth}
    \includegraphics[width=\linewidth]{dcgan/samples_epoch0010.png}
    \caption{Epoch 10}
  \end{subfigure}
  \hfill
  \begin{subfigure}[b]{0.32\linewidth}
    \includegraphics[width=\linewidth]{dcgan/samples_epoch0050.png}
    \caption{Epoch 50}
  \end{subfigure}
  \hfill
  \begin{subfigure}[b]{0.32\linewidth}
    \includegraphics[width=\linewidth]{dcgan/samples_epoch0100.png}
    \caption{Epoch 100}
  \end{subfigure}
  \caption{DCGAN sample grids at training epochs 10, 50, and 100.}
  \label{fig:dcgan_progression}
\end{figure}

\begin{figure}[H]
  \centering
  \begin{subfigure}[b]{0.32\linewidth}
    \includegraphics[width=\linewidth]{ddpm/samples_epoch0010.png}
    \caption{Epoch 10}
  \end{subfigure}
  \hfill
  \begin{subfigure}[b]{0.32\linewidth}
    \includegraphics[width=\linewidth]{ddpm/samples_epoch0100.png}
    \caption{Epoch 100}
  \end{subfigure}
  \hfill
  \begin{subfigure}[b]{0.32\linewidth}
    \includegraphics[width=\linewidth]{ddpm/samples_epoch0200.png}
    \caption{Epoch 200}
  \end{subfigure}
  \caption{DDPM sample grids at training epochs 10, 100, and 200.}
  \label{fig:ddpm_progression}
\end{figure}

\subsection{VQ-VAE Reconstruction}
\label{sec:vqvae_recon}

\begin{figure}[H]
  \centering
  \includegraphics[width=0.85\linewidth]{vqvae/recon_slide.png}
  \caption{VQ-VAE reconstructions at epoch 100 (top row: inputs, bottom row: reconstructed).}
  \label{fig:vqvae_recon}
\end{figure}

The VQ-VAE achieves a codebook perplexity of $403.3 / 512$, indicating that
$\approx 79\%$ of codes are actively used.  Reconstruction loss at convergence
is $0.0066$ (MSE), VQ commitment loss $0.0033$.  The good reconstruction
quality contrasts with the blurrier unconditional samples, which suggests that
the PixelCNN prior is a likely bottleneck: the autoencoder itself captures
sufficient detail, but the prior's ability to model the full distribution of
token sequences is limited.

\subsection{Latent Interpolation}
\label{sec:interpolation}

To assess the smoothness of the learned latent space, we interpolate between
two latent vectors $\mathbf{z}_A$ and $\mathbf{z}_B$.  For models whose prior
is $\mathcal{N}(0,I)$ (DCGAN, DDPM), we use \emph{spherical linear interpolation}
(slerp):
\begin{equation}
  \mathbf{z}_\alpha = \frac{\sin((1-\alpha)\theta)}{\sin\theta}\,\mathbf{z}_A
                    + \frac{\sin(\alpha\theta)}{\sin\theta}\,\mathbf{z}_B,
  \quad \theta = \arccos\!\left(\frac{\mathbf{z}_A \cdot \mathbf{z}_B}
                                     {\|\mathbf{z}_A\|\,\|\mathbf{z}_B\|}\right),
\end{equation}
for $\alpha \in \{0, \tfrac{1}{9}, \ldots, 1\}$.  Slerp preserves the norm of
vectors drawn from an isotropic Gaussian: linear interpolation through the
origin would traverse a low-density region of $\mathcal{N}(0,I)$ and produce
systematically blurry intermediate frames.  For VQ-VAE, we interpolate linearly
in the continuous pre-quantization encoder space, since the prior operates over
discrete token indices rather than a Gaussian.

\begin{figure}[H]
  \centering
  \begin{subfigure}[b]{\linewidth}
    \includegraphics[width=\linewidth]{dcgan/interpolation_compare.png}
    \caption{DCGAN: slerp (top row) vs.\ linear interpolation (bottom row)
             in $\mathcal{N}(0,I)$ latent space.}
  \end{subfigure}\\[8pt]
  \begin{subfigure}[b]{\linewidth}
    \includegraphics[width=\linewidth]{vqvae/interpolation.png}
    \caption{VQ-VAE: linear interpolation in continuous pre-quantization encoder space.}
  \end{subfigure}\\[8pt]
  \begin{subfigure}[b]{\linewidth}
    \includegraphics[width=\linewidth]{ddpm/interpolation_compare.png}
    \caption{DDPM: slerp (top row) vs.\ linear interpolation (bottom row)
             between initial noise tensors $x_T \sim \mathcal{N}(0,I)$.}
  \end{subfigure}
  \caption{Latent interpolation strips (10 images: 2 endpoints + 8 intermediate frames).
           For DCGAN and DDPM, both slerp and linear interpolation are shown to
           enable direct comparison.}
  \label{fig:interpolation}
\end{figure}

\paragraph{Discussion.}
Both slerp and linear interpolation are shown for DCGAN and DDPM, as the project
instructions specify linear interpolation.  The comparison illustrates why slerp
is preferable for Gaussian priors: the midpoint of a linear interpolation between
two vectors drawn from $\mathcal{N}(0,I)$ has expected norm
$\|\tfrac{1}{2}(z_0 + z_1)\| \approx \tfrac{1}{\sqrt{2}}\|z\|$, placing
intermediate latents in a low-density region near the origin.  The resulting
decoded images are noticeably blurrier in the centre of the linear strip,
particularly for DDPM where the effect of a poorly-initialised $x_T$ compounds
over 50 reverse steps.  Slerp maintains constant norm throughout the path and
avoids this artefact.  VQ-VAE uses linear interpolation in the continuous
pre-quantization encoder space, where no Gaussian shell assumption applies;
the discrete quantization step introduces occasional abrupt transitions when
token assignments change near a Voronoi cell boundary.

Smooth interpolations across all models indicate that the learned latent spaces
encode a continuous manifold rather than memorising isolated training examples.
DDPM slerp in $x_T$ space produces the most photorealistic intermediate frames.

\section{Mode Collapse Analysis}
\label{sec:mode_collapse}

It is important to distinguish between two related but distinct failure modes
observed in the DCGAN training:

\begin{description}[noitemsep]
  \item[Mode collapse] occurs when the generator maps a large region of the
        latent space $\mathcal{N}(0,I)$ to a narrow subset of the data
        distribution, producing images with limited variety.  Symptoms include
        repeated similar facial proportions and colour palettes across samples
        drawn from different $\mathbf{z}$.
  \item[Generator starvation] is an independent training-budget problem: with
        $n_{\text{critic}}=5$, the generator receives far fewer weight updates
        than the discriminator.  A poorly-trained generator will score high FID
        even without collapsing to a single mode, simply because it has not had
        enough gradient signal to learn the full distribution.
\end{description}

Both effects are present in the original WGAN-GP training run and contribute
jointly to the high FID of 221.51.  Inspecting generated grids
(Figure~\ref{fig:dcgan_progression}) confirms reduced diversity (mode-collapse
symptoms) alongside overall blurriness attributable to insufficient generator
updates.

\textbf{Mitigation strategies applied:}
\begin{itemize}[noitemsep]
  \item WGAN-GP gradient penalty ($\lambda=10$) replaces standard discriminator
        weight clipping, providing smoother gradients to the generator.
  \item Periodic sample grids saved every 10 epochs were used to monitor for
        mode collapse during training (no per-epoch FID due to cost).
  \item Note: the initial implementation contained a label smoothing parameter
        (real labels $\rightarrow 0.9$) intended for standard BCE-GAN training.
        This parameter operates on binary cross-entropy targets and has no effect
        on the WGAN-GP loss path, which uses a scalar critic score rather than
        binary labels.  It was therefore \emph{inactive} during WGAN-GP training.
\end{itemize}

The generator starvation hypothesis and corrected configurations are analysed
quantitatively in Section~\ref{sec:dcgan_bottleneck}.

\section{DCGAN Training Bottleneck Analysis}
\label{sec:dcgan_bottleneck}

\subsection{Generator Starvation}

A post-hoc analysis of the DCGAN training budget revealed a potential source of
under-training.  With $n_{\text{critic}}=5$ and 156 batches per epoch
($\lfloor 9997/64 \rfloor$ with \texttt{drop\_last=True}), the generator
receives only $\lfloor 156/5 \rfloor \times 100 = \mathbf{3\,100}$ update steps
over 100 epochs.  A standard DCGAN with $n_c=1$ would accumulate $156 \times 100
= \mathbf{15\,600}$ generator updates at the same computational cost.

\subsection{Corrected Configurations}

Two corrected configurations were trained to test whether increasing the
generator update budget improves FID:

\begin{enumerate}[noitemsep]
  \item \textbf{Standard DCGAN}: switch to standard BCE loss with $n_c=1$
        (the canonical DCGAN setting, not a compromise); increased generator
        capacity ($g_{\text{ch}}=128$); 150 epochs $\Rightarrow$
        $\mathbf{23\,400}$ G-steps.
  \item \textbf{WGAN-GP $n_c=2$}: pragmatic reduction of critic steps while
        retaining the Wasserstein loss; 200 epochs $\Rightarrow$
        $\mathbf{15\,600}$ G-steps.  Note that $n_c < 5$ weakens the
        Lipschitz guarantee; a theoretically proper fix would require
        $\approx$500 epochs.
\end{enumerate}

\begin{table}[H]
  \centering
  \caption{FID results for original and corrected DCGAN configurations.}
  \label{tab:dcgan_fixed}
  \begin{tabular}{lccc}
    \toprule
    Configuration & G-steps & FID $\downarrow$ & Notes \\
    \midrule
    Original WGAN-GP ($n_c=5$, 100 ep.) & 3\,100  & 221.51 & baseline \\
    Standard DCGAN ($n_c=1$, 150 ep.)   & 23\,400 & \textbf{98.99}  & BCE loss \\
    WGAN-GP $n_c=2$ (200 ep.)           & 15\,600 & 158.48 & Wasserstein loss \\
    \bottomrule
  \end{tabular}
\end{table}

\begin{figure}[H]
  \centering
  \begin{subfigure}[b]{0.48\linewidth}
    \includegraphics[width=\linewidth]{dcgan_fixed_vanilla/samples_epoch0150.png}
    \caption{Standard DCGAN (epoch 150)}
  \end{subfigure}
  \hfill
  \begin{subfigure}[b]{0.48\linewidth}
    \includegraphics[width=\linewidth]{dcgan_fixed_wgangp/samples_epoch0200.png}
    \caption{WGAN-GP $n_c=2$ (epoch 200)}
  \end{subfigure}
  \caption{Samples from corrected DCGAN configurations.}
  \label{fig:dcgan_fixed}
\end{figure}

\subsection{Discussion}

Both corrected configurations yield substantially lower FID.  The standard DCGAN
at 23\,400 G-steps achieves FID 98.99, a $2.2\times$ improvement over the
original.  The WGAN-GP $n_c=2$ at 15\,600 G-steps achieves FID 158.48.

These results suggest that the original WGAN-GP result was limited not only by
architecture but also by the low training budget allocated to generator updates.
However, because each corrected configuration changes several factors
simultaneously -- loss function, $n_c$, model capacity, and training duration --
these should be interpreted as diagnostic follow-up experiments rather than a
clean single-factor ablation.  The improvement from standard DCGAN over
WGAN-GP $n_c=2$ may reflect the larger generator update budget (23\,400 vs.\
15\,600 steps), the switch to BCE loss, or increased generator capacity, and
cannot be attributed to any single factor from this experiment alone.

\section{Ablation Studies}
\label{sec:ablation}

\subsection{DDPM: DDIM Inference Steps}

We evaluate the effect of the number of DDIM steps on FID without retraining
(inference-only ablation on the epoch-200 checkpoint, non-EMA weights).

\begin{table}[H]
  \centering
  \caption{DDPM FID vs.\ number of DDIM inference steps (epoch-200 checkpoint,
           non-EMA weights, 5\,000 samples).}
  \label{tab:ddpm_steps}
  \begin{tabular}{ccc}
    \toprule
    DDIM Steps & FID $\downarrow$ & Speedup vs.\ 1000-step \\
    \midrule
    10  & 66.92 & $100\times$ \\
    25  & 60.02 & $40\times$ \\
    50  & 54.51 & $20\times$ \\
    100 & 48.70 & $10\times$ \\
    200 & 39.81 & $5\times$ \\
    \bottomrule
  \end{tabular}
\end{table}

FID improves monotonically with more steps, with diminishing returns above 100.
The 50-step setting ($20\times$ speedup, FID 54.5) provides a practical
speed-quality trade-off.  Steps below 25 degrade noticeably ($+27$ FID vs.\
50 steps), making them unsuitable for qualitative evaluation.

\subsection{DDPM: EMA Weights at Inference}

Using EMA weights at inference has a large effect independent of step count.

\begin{table}[H]
  \centering
  \caption{Effect of EMA weights at 50 DDIM steps (epoch-200 checkpoint).}
  \label{tab:ddpm_ema}
  \begin{tabular}{lc}
    \toprule
    Inference weights & FID $\downarrow$ \\
    \midrule
    Without EMA & 54.51 \\
    \textbf{With EMA (decay 0.9999)} & \textbf{24.36} \\
    \bottomrule
  \end{tabular}
\end{table}

EMA weights reduce FID by 55\% at 50 steps (from 54.51 to 24.36), a larger
improvement than doubling the number of inference steps.  The EMA averages
weight noise that accumulates during training, which is particularly beneficial
for diffusion models because errors in the reverse chain compound over many steps.
The 50-step EMA configuration is therefore the default reported throughout this
work.

\subsection{VQ-VAE: Codebook Size}
\label{sec:vqvae_ablation}

The codebook size $K$ is the primary hyperparameter governing the expressiveness
of the discrete latent space.  An initial run with $K=1024$ was conducted before
settling on the final configuration.  With $K=1024$, codebook perplexity
collapsed to $\approx 300/1024$ (29\% utilisation) during training, meaning
over 70\% of codes were never assigned.  Switching to $K=512$ with EMA codebook
updates (decay $\gamma=0.99$) recovered utilisation to 79\% ($403/512$) and
produced visibly sharper reconstructions.  Full multi-seed ablation for VQ-VAE
was not conducted due to the 200-epoch (Phase 1 + Phase 2) training cost per run.
Other hyperparameters (embedding dimension, commitment cost $\beta$, PixelCNN
depth) were set based on the original VQ-VAE paper \cite{vqvae} and not
systematically swept.

\subsection{DCGAN: Latent Dimension and Learning Rate}
\label{sec:dcgan_ablation}

We train DCGAN configurations for 100 epochs each, varying the latent
dimension $z_{\text{dim}}$ and learning rate independently across three seeds
per configuration (all other hyperparameters fixed to baseline:
$z_{\text{dim}}=128$, lr $=2\times10^{-4}$, $n_c=5$).  FID is reported as
mean $\pm$ population standard deviation across seeds.

\begin{table}[H]
  \centering
  \caption{DCGAN ablation: FID (mean $\pm$ std across 3 seeds, 100 epochs).}
  \label{tab:dcgan_ablation}
  \begin{tabular}{lcc}
    \toprule
    Configuration & FID mean $\pm$ std & $\Delta$ vs.\ baseline \\
    \midrule
    $z_{\text{dim}} = 64$              & $246.65 \pm 5.57$  & $-17.60$ \\
    $z_{\text{dim}} = 128$ (baseline)  & $264.25 \pm 20.21$ & -- \\
    $z_{\text{dim}} = 256$             & $266.05 \pm 2.52$  & $+1.80$ \\
    \midrule
    lr $= 1 \times 10^{-4}$            & $295.11 \pm 18.27$ & $+30.86$ \\
    lr $= 2 \times 10^{-4}$ (baseline) & $264.25 \pm 20.21$ & -- \\
    lr $= 4 \times 10^{-4}$            & $220.05 \pm 4.47$  & $-44.20$ \\
    \midrule
    Augmented ($z_{128}$, lr $2\times10^{-4}$) & $284.05 \pm 14.80$ & $+19.80$ \\
    \bottomrule
  \end{tabular}
\end{table}

Table~\ref{tab:dcgan_raw} lists the individual per-seed FID values underlying the
summary statistics.

\begin{table}[H]
  \centering
  \caption{DCGAN ablation: raw FID per seed (seeds 0, 42, 123; 100 epochs each).}
  \label{tab:dcgan_raw}
  \small
  \begin{tabular}{lrrrcc}
    \toprule
    Configuration & Seed 0 & Seed 42 & Seed 123 & Mean & Std \\
    \midrule
    $z_{\text{dim}}=64$              & 254.40 & 241.55 & 244.01 & 246.65 & 5.57 \\
    $z_{\text{dim}}=128$ (baseline)  & 239.25 & 288.73 & 264.78 & 264.25 & 20.21 \\
    $z_{\text{dim}}=256$             & 264.20 & 264.33 & 269.61 & 266.05 & 2.52 \\
    \midrule
    lr $= 1\times10^{-4}$            & 313.98 & 270.39 & 300.94 & 295.11 & 18.27 \\
    lr $= 2\times10^{-4}$ (baseline) & 239.25 & 288.73 & 264.78 & 264.25 & 20.21 \\
    lr $= 4\times10^{-4}$            & 226.01 & 215.26 & 218.87 & 220.05 & 4.47 \\
    \midrule
    Augmented                         & 304.45 & 277.92 & 269.78 & 284.05 & 14.80 \\
    \bottomrule
  \end{tabular}
\end{table}

Latent dimension has a small effect: the differences between
$z_{\text{dim}} \in \{64, 128, 256\}$ are within the baseline seed variance
($\pm20$), so no firm conclusion about the optimal $z_{\text{dim}}$ can be
drawn.  The high baseline variance itself is notable and indicates that DCGAN
training outcomes are sensitive to random initialisation.

Learning rate has a much larger impact.  A higher lr of $4\times10^{-4}$
reduces mean FID by 44 points with low variance ($\pm4.47$), while the lowest
lr ($1\times10^{-4}$) increases mean FID by 31 points.  The likely mechanism is
that a low learning rate causes the generator to update too slowly to keep pace
with the well-trained critic, resulting in poor gradient signal.

Enabling data augmentation hurts FID ($284.05 \pm 14.80$ vs.\ $264.25 \pm 20.21$
baseline).  In this experiment, augmentation was applied only to the real images
via the DataLoader, not to the generated images before they were passed to the
discriminator.  As a result, the discriminator was trained to distinguish
augmented real images (labelled real) from clean generated images (labelled fake),
causing it to accept distortions as a signature of real data and
incorrectly steer the generator toward augmentation artefacts rather than the
true data distribution.  The correct approach -- applying identical transforms to
both real and generated images before the discriminator -- is implemented in
strategies such as ADA \cite{karras2020training}, but was not applied here.
This is a known pitfall in GAN augmentation and limits the conclusions that can
be drawn from this ablation.

\section{Cats and Dogs Experiment}
\label{sec:cats_and_dogs}

\subsection{Setup}
The Kaggle dogs-vs-cats dataset \cite{dogsvscats} provides 12\,500 cat images
and 12\,500 dog images (25\,000 total).  We train the same DCGAN architecture
(100 epochs, all hyperparameters identical to the cats-only baseline) on the
combined dataset without class labels, as an unconditional generative model.

DCGAN was selected for this experiment rather than DDPM for two reasons.
First, training speed: DCGAN trains in $\sim$2\,h vs.\ $\sim$12\,h for DDPM,
making it feasible to run an exploratory 100-epoch experiment within the
available compute budget.  Second, the experiment addresses a dataset-composition
question (does adding dog images help?) rather than a quality benchmark; the
relative change in FID is the key signal, not the absolute value.  A
higher-quality investigation would replicate this experiment with DDPM, but
that would require $\approx$30\,h of additional compute on the 25\,000-image
dataset.

\subsection{Evaluation}
FID is computed separately against 5\,000 real cat images and 5\,000 real dog
images to measure how well the generated distribution covers each class.

\begin{table}[H]
  \centering
  \caption{DCGAN (cats+dogs) FID compared to cats-only model.}
  \label{tab:cats_and_dogs}
  \begin{tabular}{lccc}
    \toprule
    Model & FID vs.\ Cats $\downarrow$ & FID vs.\ Dogs $\downarrow$ & FID vs.\ Combined $\downarrow$ \\
    \midrule
    DCGAN (cats only, 100 ep.) & 221.51 & -- & -- \\
    DCGAN (cats+dogs, 100 ep.) & 206.57 & 197.79 & \textbf{190.39} \\
    \bottomrule
  \end{tabular}
\end{table}

\begin{figure}[H]
  \centering
  \begin{subfigure}[b]{0.48\linewidth}
    \includegraphics[width=\linewidth]{cats_and_dogs/samples_epoch0025.png}
    \caption{Epoch 25}
  \end{subfigure}
  \hfill
  \begin{subfigure}[b]{0.48\linewidth}
    \includegraphics[width=\linewidth]{cats_and_dogs/samples_epoch0100.png}
    \caption{Epoch 100}
  \end{subfigure}
  \caption{DCGAN (cats+dogs) sample grids at epochs 25 and 100.}
  \label{fig:cats_dogs_samples}
\end{figure}

\subsection{Qualitative Assessment}

Inspecting the generated samples (Figure~\ref{fig:cats_dogs_samples}), the model
produces a mixture of recognisable and ambiguous outputs.  Some images display
distinctly cat-like features: rounded faces, pointed ears, and compact body
proportions.  Others show dog-like characteristics: longer snouts, floppy ears,
and leaner body shapes.  However, a large proportion of samples are visually
ambiguous -- dark quadruped silhouettes without clearly discernible species
identity.  We observe no systematic collapse to one class: both cat-like and
dog-like structures are present throughout the grid.  This suggests the
unconditional model has learned a blended distribution that captures shared
low-level animal features (fur texture, eyes, body outline) while sometimes
failing to commit to either species -- a predictable consequence of training
without class labels on a mixed dataset.  A classifier-based analysis (e.g.,
using a pre-trained cat/dog classifier on the generated samples) would be needed
for a quantitative confirmation of this qualitative observation.

\subsection{Quantitative Discussion}
Training an unconditional generative model on a two-class dataset poses an
interesting challenge: without conditioning, the model cannot guarantee
class-distinct generation.

The cats+dogs DCGAN achieves FID $206.57$ against real cats, 15 points lower
than the cats-only model's FID of $221.51$.  This improvement is suggestive of
a beneficial effect: additional dog images may expose the model to a wider
variety of animal textures and structures, with shared low-level features (fur,
eyes, snout shape) acting as implicit data augmentation.  However, this
interpretation should be treated cautiously.  The 15-point improvement is
comparable to the seed variance observed in the DCGAN ablation
($\sigma \approx 20$ FID), so the result is not statistically conclusive from a
single run.  Furthermore, FID alone cannot confirm whether the generated images
are visually distinct as cats or dogs versus blended representations; a
classifier-based analysis would be needed to verify class identity in generated
samples.

The FID against real dogs ($197.79$) is similar to the cat FID ($206.57$),
suggesting the model does not strongly collapse to one class.

As a more direct measure of how well the generated distribution covers the
training distribution, we also computed FID against a combined reference set of
2\,500 real cats and 2\,500 real dogs ($N=5\,000$ total).  The resulting FID is
$190.39$, lower than either individual class FID.  This is the most
methodologically appropriate comparison for an unconditional model trained on a
mixed dataset: the reference set reflects the same class balance as the training
data.  The improvement over the per-class FIDs (from $\approx202$ average to
$190$) is consistent with the model covering a blended distribution rather than
cleanly separating into cat and dog modes.

\section{Conclusion}
\label{sec:conclusion}

We trained and evaluated three families of generative models on cat images:

\begin{itemize}[noitemsep]
  \item \textbf{DDPM} achieves the best FID (24.36 with EMA at 50 DDIM steps)
        and the most visually coherent samples.  The 50-step DDIM configuration
        provides a practical speed-quality trade-off ($20\times$ speedup over
        full sampling).  EMA weights are the single most impactful component:
        they reduce FID by 55\% relative to non-EMA inference at the same
        step count.
  \item \textbf{VQ-VAE + PixelCNN} achieves FID 187.54.  Reconstruction
        quality is good (MSE $= 0.0066$, codebook utilisation 79\%), but the gap
        between reconstruction quality and sample quality suggests the PixelCNN
        prior is a likely bottleneck for unconditional generation.
  \item \textbf{DCGAN (WGAN-GP)} achieves FID 221.51, with visible mode
        collapse symptoms despite WGAN-GP regularisation.  Training is the
        fastest ($\sim$2\,h).  Post-hoc analysis suggests that generator
        starvation ($n_c=5$ yielding only 3\,100 G-steps) contributed to the
        poor result; corrected configurations reduce FID to 98.99 (standard
        DCGAN, 23\,400 G-steps) and 158.48 (WGAN-GP $n_c=2$, 15\,600 G-steps).
\end{itemize}

These results are consistent with the literature showing that diffusion models
offer superior image quality compared to GAN and VAE approaches on
small-scale image datasets with limited compute \cite{diffusionbeatsgan}.

\paragraph{Methodological limitations.}
The main FID comparisons are single-run results with no standard deviation
reported for DDPM and VQ-VAE, because each additional run costs $\sim$12\,h
(DDPM) or $\sim$4\,h (VQ-VAE).  The DCGAN ablation in Section~\ref{sec:dcgan_ablation}
was repeated across three seeds to quantify variance; ideally the same rigour
would apply to all models.  The relative ranking between architectures is
reliable for large FID gaps (DDPM vs.\ the rest) but the VQ-VAE vs.\ DCGAN
difference ($\approx34$ FID, or $\approx100$ FID after the DCGAN fix) may not
be fully robust to seed variation.

The DDPM ablation covers inference-time parameters only (DDIM steps and EMA).
Training-time hyperparameters -- noise schedule type, number of diffusion steps
$T$, U-Net channel multipliers, and learning rate -- were not systematically
ablated; a full grid search would require multiple retrains.

The comparison is not compute-normalised.  No held-out test set was used: FID
reflects similarity to the training distribution.  The cats-and-dogs improvement
and the augmentation penalty are both within the observed DCGAN seed variance
($\approx \pm 20$ FID) and should therefore be treated as exploratory findings.
The corrected DCGAN experiments change multiple factors simultaneously, so the
improvement cannot be attributed to generator starvation alone.

Augmentation in the GAN ablation was applied only to real images via the
DataLoader, not to generated images before the discriminator -- the correct
approach for GANs (as implemented in ADA \cite{karras2020training}) would apply
identical transforms to both.  This limits the conclusions of the augmentation
ablation.

\textbf{DDIM timestep schedule.}
The DDIM sampling implementation uses \texttt{range(0, T, T//S)} reversed,
producing a grid $\{S\cdot\lfloor T/S\rfloor,\ldots,\lfloor T/S\rfloor, 0\}$
starting from $t=980$ rather than $t=999$ for $T=1000$, $S=50$.
Since $\bar\alpha_{980} \approx 0.002$ (near pure noise), the practical impact
on sample quality is negligible, but the schedule does not strictly begin at the
final diffusion step.

\textbf{FID reference set.}
Real reference images for FID were saved once using a sequential DataLoader
(no shuffle) and reused across all evaluations.  Since the same fixed set of
5\,000 images was used for every model comparison, the FID values are
internally consistent.  The images are drawn from the first 5\,000 samples
in file-sort order across all subdirectories; diversity across subdirectories
CAT\_00--CAT\_06 was not explicitly verified but the dataset is known to
contain a broad range of cat breeds and poses.  FID sample directories were
not cleared between runs, but sequential file naming (000000.png, \ldots)
means existing files are overwritten rather than accumulated.

\paragraph{Potential improvements and future work.}
\begin{itemize}[noitemsep]
  \item Training DDPM at $128\times128$ resolution with gradient checkpointing
        would likely improve FID further.
  \item A class-conditional DCGAN (e.g., projection discriminator) trained on
        the cats-and-dogs dataset would enable targeted cat/dog generation and
        allow a cleaner assessment of multi-class transfer.
  \item Replacing the PixelCNN prior with a Transformer-based prior (e.g.,
        VQ-VAE-2 \cite{vqvae}) could improve sample diversity.
  \item A cosine noise schedule for DDPM may improve results at lower resolutions.
\end{itemize}

\section*{AI Usage Statement}
\label{sec:ai}

GitHub Copilot (Claude Sonnet) was used as a coding assistant throughout this
project for: boilerplate generation, debugging PyTorch errors, and suggestions
for hyperparameter configurations.  All architectural decisions, experimental
design, and analysis were performed by the author.  Generated code was reviewed
and adapted before use.

\section*{Reproducibility}

All code is available in the project repository.  To reproduce results:

\begin{enumerate}[noitemsep]
  \item Install dependencies: \texttt{pip install -r requirements.txt}
  \item Download and place the CAT Dataset in \texttt{data/cats/}
  \item Train models: \texttt{python -m src.training.train\_dcgan --config configs/dcgan\_config.yaml}, analogously for VQ-VAE and DDPM.
  \item Evaluate: \texttt{python evaluate\_all.py}
  \item Run ablations: \texttt{python ablate\_ddpm\_steps.py} and \texttt{python ablate\_dcgan\_multiseed.py}
\end{enumerate}

Global seed 42 is set via \texttt{src/utils/seed.py} for full reproducibility.

\bibliographystyle{plain}
\bibliography{references}

\end{document}
